% Addition TER-01--04, TER-06--07. TER-05 remains developed in 09_area,
% 09b_hojas_memoria, 09c_entropia_sectorial and the subsequent composition.
% Sources: BOLTZMANN_DESARROLLO_OPERATORIO_20260927/{retorno,normalizacion,
% incidencia,seleccion,accion_horizonte_20260927}; FOCK_PARIDAD_RAMAS_Y_SELECTOR;
% REVISION_LIFT_POSITIVO_Y_ESPECTRO; ADDENDUM_TRANSPORTE_TERMICO_Y_SECCIONES.
% Additive edition: the formal kernel and earlier proofs are preserved.
% Inherited causal receipt: APP retains sum, product, remainders and quotients.
% TRIT retains regime and orientation in the received calendar.
% TPK transports path, carry, boundary and memory in the nonadic connection.
% The enriched state prolongs the record when phase returns.
% The joint discrete structure of the continuum is the inherited domain.
% Examined HMT output: capacity, graded carrier and marked reader.
% The conventional spectral realization is subsequent and retains declared rules.
\section{Spectral capacity, energy-bearing memory and thermal transport}
\label{sec:ii-ter27-antecedentes}

The following construction combines extension capacity, renewal of the two
sheets and their spectral realizations. Its purpose is to determine what a
thermal reading preserves when phase returns while the record continues to
grow. Its operations originate in the enriched APP--TRIT--TPK state: the
signature catalogue, calendar, orientation and memory precede the traces
to be evaluated. The linear spaces, compressions and preparations added
to those antecedents are defined explicitly. Each result thus retains its
construction conditions and its role in the argument.

\subsection{Capacity, vacancy and joint renewal}
\label{sec:ii-ter27-capacidad}

The comparison between the $729$ six-trit words and the $1000$ positions
of a decimal block defines, with the origin fixed in
\ref{vac-capacidad-trit},
\begin{equation}
 C_t=\lfloor(t+1)\rho\rfloor,
 \qquad \rho=\log_{1000}729=\log_{10}9,
 \qquad z_t=1-(C_t-C_{t-1}).
 \label{eq:ii-ter27-capacidad}
\end{equation}
Since $0<\rho<1$, one has $z_t\in\{0,1\}$. Capacity counts blocks in
this comparison; its exact evaluation uses
$10^{C_t}\leq9^{t+1}<10^{C_t+1}$. For $1\leq t\leq20$, $C_t=t$,
whereas $C_{21}=20$. Indeed, for $t\leq20$ the lower inequality is
equivalent to $9(9/10)^t\geq1$, and checking $t=20$ suffices; at $t=21$
it reverses. The remaining comparisons place $9^{22}$ between
$10^{20}$ and $10^{21}$. The first positive vacancy is therefore
$p_1=21$, with
\[
 C_{20}=C_{21}=20,\qquad C_{22}=21,\quad C_{23}=22,\quad C_{24}=23.
\]
The twenty capacity increments provide twenty effective labels. The
cut index $24$, the $25$ positions from $0$ to $24$, and capacity $23$
are records of different types.

The joint emission catalogue is
$\Omega=\mathcal S_+\times\mathcal S_\times$, with cardinalities
$18$ and $26$. On $\ell^2(\Omega)$ with uniform counting measure,
renewal of the active coordinate is represented by
\[
 E_+=\frac{\mathbf J_{18}}{18}\otimes I_{26},\qquad
 E_\times=I_{18}\otimes\frac{\mathbf J_{26}}{26}.
\]
These are commuting orthogonal projections. Their product is the
constant projection $\Pi$, all of whose matrix entries equal $1/468$.
The reduced phase reading uses the departure regime:
\[
 K_{\rm ph}((u,r),(v,r+1))=P_{\tau(r)}(u,v),
 \qquad(P_+,P_-,P_0)=(E_+,E_\times,I),
 \quad\tau(r)=M_{\rm ph}(1+r).
\]
This operation follows the enriched dynamics and its holonomy, in the
order $U_t\longrightarrow\Gamma_9\longrightarrow K_{\rm ph}$.

\begin{proposition}[Two-sheet renewal on the capacity clock]
\label{prop:ii-ter27-renovacion}
Evaluate the preceding reader through $\mathcal R_t=K_{\rm ph}^{a_t}$,
where $a_t=C_{t+1}-C_t$. Starting at the first vacancy, the earliest
instant at which its matrix output is independent of the initial emission
is $24$. The successive ranks are $468,26,1$; in the conjugate
orientation they are $468,18,1$.
\end{proposition}
\begin{proof}
The departure capacities $20,21,22$ correspond to regimes $0,+1,-1$.
The transitions $21\to22\to23\to24$ therefore apply
$I,E_+,E_\times$, with accumulated products $I,E_+,\Pi$.
Tensor-product ranks multiply: $\operatorname{rank}E_+=26$,
$\operatorname{rank}E_\times=18$, and $\operatorname{rank}\Pi=1$.
The rows of $I$ and $E_+$ depend on the initial emission; the rows of
$\Pi$ coincide. This establishes minimality. The conjugate sequence
applies $I,E_\times,E_+$ and has the same final product. Composition
through capacity preserves chronological time:
\[
 (t,C_t,v_t)\mapsto(t+1,C_t+a_t,v_t+1-a_t),\qquad v_t=t-C_t.
\]
The identity $t=C_t+v_t$ is maintained also when $a_t=0$.
\end{proof}

Reduction from phase modulo nine to regime modulo three is exact for
this reader. With $Q(u,r)=(u,[r]_3)$,
\[
 \sum_{s:[s]_3=b}K_{\rm ph}((u,r),(v,s))
 =\mathbf1_{b=[r+1]_3}P_{\tau(r)}(u,v).
\]
The sum depends only on $Q(u,r)$, proving strong lumpability. The lift,
however, retains $C=9m+3b+a$. The transition $20\to23$ sends
$(m,b,a)=(2,0,2)$ to $(2,1,2)$: $a$ returns and $b$ advances. The
positive phase modulo nine changes from $3$ to $6$. Uniform renewal
is a rank-one matrix output; the earlier microstate and path remain in
the enriched archive rather than in that matrix output alone.
The information of a uniform regime is $\log3$; that of renewed
emissions is $\log468$; the entropy rate of a deterministic rotation
is zero. These readings refer to different ensembles and operations.

\subsection{Capacity returns and cylinder valuation}
\label{sec:ii-ter27-cociclo}

Define strict regime return by
\[
 R_3(s)=\min\{t>s:C_t>C_s,\ C_t\equiv C_s\pmod3\}.
\]
The sequence $C_t$ is unbounded and has increments zero or one.
It therefore reaches $C_s+3$ before any larger positive value congruent
to $C_s$, so $C_{R_3(s)}=C_s+3$. Starting at $s_0=24$, the corrected
localization is
\begin{equation}
 s_j=\min\{t\geq24:C_t=23+3j\},\qquad C_{s_j}=23+3j.
 \label{eq:ii-ter27-localizacion}
\end{equation}
The first times are $24,27,30$. Later intervals may contain vacancies;
the capacity increase nevertheless remains exactly three.

To explain the valuation, write each block as
$d_r=100a_{r,1}+10a_{r,2}+a_{r,3}$, with $0\leq a_{r,k}\leq9$.
A prefix of $c$ blocks is transported to the three integers
$X_k=\sum_{r=1}^c a_{r,k}10^{c-r}$. Euclidean division gives a
bijection; its inverse regroups the digits at each position. Truncating
one block truncates one digit in each coordinate. Addition with carry
is transported through this complete bijection, retaining its quotients.

The scalar cylinder of length $1000^{-c}$ corresponds to the half-open cube
\[
 \prod_{k=1}^3[X_k10^{-c},(X_k+1)10^{-c}),\qquad
 J_c=10^{-c}I_3,\qquad\det J_c=1000^{-c}.
\]
The equality expresses a correspondence of cylinder measures. Between
two capacities, $J_{c',c}=10^{-(c'-c)}I_3$ satisfies
$J_{c'',c'}J_{c',c}=J_{c'',c}$. The positive isotropic similarity is
determined by its determinant, because $s^3=1000^{-(c'-c)}$ has one
positive root. At the marked returns,
\begin{equation}
 J_{C_{s_j}}=10^{-23}1000^{-j}I_3.
 \label{eq:ii-ter27-cociclo}
\end{equation}
The factor $10^{-23}$ is the scale of one cube coordinate; the scalar
cylinder measure at that capacity is $10^{-69}$. This distinction
preserves the geometric degree of the reader. An extractor fixed at
$s_0$, followed by compatible truncation, satisfies
$E_n\pi_{m,n}=E_m$ for $m\geq n\geq s_0$. Extension thus preserves
the marked record without recomputing its ranks from a later cut.

\subsection{Relation carriers and their interchange realization}
\label{sec:ii-ter27-portador}

Let $U_m$ be the space of capacity labels and $P_p$ the space of prefix
positions. The diagonal expectation on $\operatorname{End}(U_m)$ is
$E_{\rm diag}(X)=\sum_iE_{ii}XE_{ii}$. Its orthogonal complement has
basis $E_{ij}$, $i\ne j$, and dimension $m(m-1)$. The second record
space is $\operatorname{End}(P_p)\oplus A_{p-1}$, where $A_{p-1}$
retains adjacent transitions as typed observations. Its dimension is
$p^2+p-1$. For $m=20$, $p=25$, the ranks are $380$ and $649$.
This choice of spaces and cut defines a realization of the reader;
the identities below prove its operations without attributing physical
selection of that domain to the count alone.

Based reversal acts by $E_{ij}\mapsto E_{m-1-i,m-1-j}$ and, on the
second space, by $E_{ij}\mapsto E_{p-1-i,p-1-j}$ and
$a_j\mapsto a_{p-2-j}$. For $m=20$ there are $190$ two-element
orbits. For $p=25$ there is exactly one fixed matrix element,
$E_{12,12}$; the remaining elements form $312$ matrix orbits and
$12$ transition orbits. Thus there is one fixed line and $324$
regular representations of $C_2$, with eigenspace dimensions $325$
and $324$. This is the antecedent of $649=1+2\cdot324$; it preserves
record reversal rather than specifying an energy grading.

A second realization uses the same interchange $J$ of two orientations
and retains mixed tensor terms. Let $W,V$ have dimensions $m,n$.
In the exterior sector of $W_+\oplus W_-$ and the symmetric oscillators
of $V_+\oplus V_-$, the fixed characters are
\begin{align}
 \chi_F(q)&=\tfrac12\bigl((1+q)^{2m}+(1-q^2)^m\bigr),\\
 \chi_B(q)&=\tfrac12\bigl(P_n(q)^2+P_n(q^2)\bigr),
 &P_n(q)&=\prod_{r\geq1}(1-q^r)^{-n}.
 \label{eq:ii-ter27-caracteres}
\end{align}
\begin{proof}
On one particle, the eigenspaces $J=\pm1$ each have dimension $m$.
The inserted trace on the exterior algebra is
$(1+q)^m(1-q)^m=(1-q^2)^m$. The projection $(I+J)/2$ averages this
trace with the ordinary character. Its degree-two coefficient is
$\bigl(\binom{2m}{2}-m\bigr)/2=m(m-1)$.
At each bosonic frequency, both signs have multiplicity $n$, giving
inserted trace $\prod_r[(1-q^r)(1+q^r)]^{-n}=P_n(q^2)$.
The ordinary coefficients in degrees zero, one and two are
$1,2n,2n^2+3n$, whereas the inserted coefficients are $1,0,n$.
Their averages are $1,n,n(n+2)$. Every formal coefficient depends
on finitely many factors, so no further analytic convergence
assumption is required.
\end{proof}
The additional $n$ in energy two comes from frequency-two creation
operators. Indeed,
$M(V)_2=V(-2)\oplus\operatorname{Sym}^2(V(-1))$; at rank $24$ its
dimension is $24+300=324$. For the two interchanged orientations,
the bosonic fixed space accumulated through energy two has dimension
$1+n+n(n+2)=(n+1)^2+n$, hence $649$ for $n=24$.
Linear identification with the earlier matrix-and-transition space
requires transport of its bases; equal dimensions do not automatically
identify Hamiltonians. The homogeneous exterior cut and accumulated
bosonic cut remain declared components of this reader.

\subsection{Multiplicative seeds and exact fibre compensation}
\label{sec:ii-ter27-fibras}

The multiplicative seed support is
$X_\times=(\mathbb Z/9\mathbb Z)^2\times\{N,E,S,O\}$, with
cardinality $324$. The six-window emitter
$e:X_\times\to\mathcal S_\times$ of the formal kernel produces
$26$ signatures. Its fibres have size $8$ for $24$ signatures,
$24$ for $555555$, and $108$ for $000000$, summing to $324$.
The temporal set $I=\{1,\ldots,24\}$ retains a type distinct from
the fibre of $555555$, despite their equal cardinalities.

Construct the microscopic space
\[
 \widetilde T=\{\Omega\}\sqcup\{a_i:i\in I\}
 \sqcup\{b_{i,y}:i\in I,y\in X_\times\},\qquad|\widetilde T|=7801.
\]
The reading $p(b_{i,y})=b_{i,e(y)}$, fixing the other records, has
image $T$ of size $1+24+24\cdot26=649$. Put $m_e=|e^{-1}(e)|$.
Assign mass $1$ to $\Omega,a_i$ and $1/m_{e(y)}$ to $b_{i,y}$.
This defines a measure $\nu$ whose fibres under $p$ all have mass one.
For $r>0$, the rates are $r$ in both directions between $\Omega$ and
$a_i$, and $r/m_{e(y)}$ from $a_i$ to $b_{i,y}$, with reverse rate $r$.

\begin{proposition}[Compensated transport of the signature tree]
\label{prop:ii-ter27-arbol}
The microscopic generator is reversible for $\nu$. The map $p$ induces
the symmetric tree on $T$, and the lift $J_pg=g\circ p$ is an isometry
intertwining the two generators.
\end{proposition}
\begin{proof}
Each terminal edge satisfies
$\nu(a_i)r/m_{e(y)}=\nu(b_{i,y})r$. Moreover,
$\sum_{y:e(y)=e}r/m_e=r$, while the reverse rate has the same value
from every representative. This proves strong lumpability.
The identity $\sum_{x:p(x)=z}\nu(x)=1$ proves $J_p^*J_p=I$;
the rate sums prove $\widetilde L J_p=J_pL_T$. The tree is connected
and has $648$ edges, hence its normalized stationary measure is uniform.
\end{proof}
The cursor half-turn induces a signature involution preserving $m_e$.
Together with $i\mapsto25-i$, it preserves $\nu$ and the rates; the
intertwiner also transports this reversal. Seeds remain recorded in
$\widetilde T$ when $T$ evaluates only their signatures.
Selection of a mark and compensated signature renewal are the operations
of this reader. Their domain and rates remain explicit: cursor paths
and the full enriched update retain additional operations. The tree
connects the internal bosonic levels; communication between all three
grades is constructed in \ref{sec:ii-ter27-corrientes}.

\subsection{Double grading and recovery from its marginals}
\label{sec:ii-ter27-doble}

The carrier selected for comparison is
\[
 \mathcal D=\mathbb C\Omega_{\rm ext}\oplus F_2^J\oplus B_{\leq2}^J.
\]
The operator $L$ retains Fock degree; $N$ assigns external grades
$0,1,2$ to the three summands. For $m=20,n=24$, their joint
decomposition is
\begin{equation}
 (N,L;d)=(0,0;1),(1,2;380),(2,0;1),(2,1;24),(2,2;624).
 \label{eq:ii-ter27-celdas}
\end{equation}
The two vacua are orthogonal and retain their provenance. The bivariate
function and its specializations are
\begin{align}
 \mathcal Z(u,v)&=1+380uv^2+u^2(1+24v+624v^2),\\
 Z_N(q)&=\mathcal Z(q,1)=1+380q+649q^2,\\
 Z_L(q)&=\mathcal Z(1,q)=2+24q+1004q^2.
 \label{eq:ii-ter27-particiones}
\end{align}
Each monomial follows by applying $u^Nv^L$ on a cell of
\eqref{eq:ii-ter27-celdas} and taking its trace. At $q=1/1000$ the
values are respectively $1.380649$ and $2.025004$. The latter contains
two vacua, $24$ energy-one states and $1004$ energy-two states.

The complete marginals $n_i,l_j$ of a distribution on these five cells
recover the cell masses:
\[
 p_{00}=n_0,\quad p_{12}=n_1,\quad p_{21}=l_1,
 \quad p_{20}=l_0-n_0,\quad p_{22}=l_2-n_1.
\]
Summation along rows and columns proves recovery when these differences
are nonnegative and the totals agree. The support is a bipartite tree;
recovery concerns its five masses, while coherences and internal indices
remain in the full state. Two isolated partition values do not replace
those marginals. As a control, $Z_N-Z_L=(q-1)(1-355q)$: at $q=1/355$
the traces coincide, but the internal-vacuum probabilities are
$1/261574$ and $126025/261574$. Equal scalars coexist with distinct states.

The based reversal of the prefix reader also determines the orientation
information omitted by its quotient. There are $190$ pairs in degree
one and $324$ in degree two; each equiprobable pair contributes $\log2$
times its mass. Hence
\[
 I_{\rm or}(q)=\frac{380q+648q^2}{Z_N(q)}\log2.
\]
Restoring the orientation label recovers that information. This quotient
and its entropy reading remain distinct from a physical erasure operation.

\subsection{Positive energy dilation of memory}
\label{sec:ii-ter27-memoria}

Since $N,L$ commute, $Q_c=N-L$ takes values $0,-1,2,1,0$ on the
preceding cells. Let $J_m=\operatorname{diag}(0,1,2,3)$ and $g>0$.
Define the isometry
\begin{equation}
 V|n,l,i\rangle=|n,l,i\rangle\otimes|n-l+1\rangle,
 \quad H_F=gL,\quad H_m=gJ_m.
 \label{eq:ii-ter27-lift}
\end{equation}
The auxiliary levels are $1,0,3,2,1$. The first coordinate retains
the complete basis, so $V^*V=I$. On each vector,
\begin{equation}
 (H_F\otimes I+I\otimes H_m)V=Vg(N+I),
 \label{eq:ii-ter27-energia-memoria}
\end{equation}
because $l+(n-l+1)=n+1$. Auxiliary energy is nonnegative and total
energies are $g,2g,3g,3g,3g$. The image is reducing for this diagonal
Hamiltonian. For $q=e^{-\beta g}$, functional calculus gives
\[
 V^*e^{-\beta H_{\rm tot}}V=q^{N+I},\qquad
 \frac{V^*e^{-\beta H_{\rm tot}}V}
 {\operatorname{Tr}(V^*e^{-\beta H_{\rm tot}}V)}=\frac{q^N}{Z_N(q)}.
\]
The filter $M_q=q^{(Q_c+I)/2}$ satisfies
$V^*(I\otimes q^{J_m})V=M_q^*M_q$ and
$q^L M_q^*M_q=q^{N+I}$. For $0<q<1$, it is contractive and
conditionally transforms $\rho_L=q^L/Z_L$ into $\rho_N=q^N/Z_N$
with probability $qZ_N/Z_L$. Recording its outcome also retains the
complementary outcome. The isometry $V$ preserves coherences; partial
trace over memory is a different operation, removing precisely the
coherences between different values of $Q_c$.
The shift $+1$ is the least integer making every auxiliary level
nonnegative. A larger common shift preserves the normalized state and
changes the recorded energy origin.

\subsection{Memory loop and decimal spectrum}
\label{sec:ii-ter27-espectro}

The Gaussian realization receives the collective $8:1$ decomposition
of nine copies and the unit memory loop. Prepare both modes in the centred
pure Gaussian reference of joint covariance $I\oplus I$, in the same
symplectic coordinates used for the following matrices:
\[
 Q_8=\frac13\begin{pmatrix}\sqrt8 I&-I\\I&\sqrt8 I\end{pmatrix},
 \quad H_1=\begin{pmatrix}2&1\\1&1\end{pmatrix},
 \quad U_{H_1}=Q_8^T\operatorname{diag}(I,H_1)Q_8.
\]
These matrices inherit their preparation and loop from system--memory
transport; the thermal calculation follows. The visible block of
$U_{H_1}U_{H_1}^T$ is
\[
 W=\frac{8I+H_1H_1^T}{9}
 =\frac19\begin{pmatrix}13&3\\3&10\end{pmatrix},
 \qquad\det W=\frac{121}{81},\quad\nu=\frac{11}{9}.
\]
In dimension two, $S=\sqrt\nu W^{-1/2}$ has determinant one and is
symplectic, with $SWS^T=\nu I$. Its metaplectic realization reduces
the Gaussian state to that isotropic covariance. The geometric state
$(1-r)\sum_{j\geq0}r^j|j\rangle\langle j|$ has mean occupation
$r/(1-r)$ and covariance $(1+r)/(1-r)I$. Uniqueness of a centred
Gaussian state with given covariance implies
\begin{equation}
 r=\frac{\nu-1}{\nu+1}=\frac1{10},\qquad
 \rho_r=(1-r)\sum_{j\geq0}r^j|j\rangle\langle j|.
 \label{eq:ii-ter27-williamson}
\end{equation}
The decimal ratio follows from the stated loop and preparation. For
example, $H_0=I$ would give $\nu=1$ and $r=0$, displaying the result's
effective dependence on the memory operator.

Euclidean division $j=3w+b$, $0\leq b<3$, defines a unitary
$U_3|j\rangle=|w,b\rangle$, with inverse
$|w,b\rangle\mapsto|3w+b\rangle$. The factorization
$1-r^3=(1-r)(1+r+r^2)$ proves, term by term,
\begin{equation}
 U_3\rho_rU_3^*=\rho_q\otimes\sigma_r,
 \quad q=r^3=\frac1{1000},\quad
 \sigma_r=\frac{\operatorname{diag}(100,10,1)}{111}.
 \label{eq:ii-ter27-agrupamiento}
\end{equation}
The unilateral shift is transported to
\[
 T=I\otimes(|1\rangle\langle0|+|2\rangle\langle1|)
       +S_w\otimes|0\rangle\langle2|,\qquad T^3=S_w\otimes I_3.
\]
Applying this formula to each basis vector proves the identity and
exhibits the carry from remainder to quotient. One has $T^*T=I$ and
$TT^*=I-|0,0\rangle\langle0,0|$; remainder return retains the
occupation advance. Acting as $U_3\otimes I$, the unitary also
transports every purification of the enlarged system.

For the energy section $H_W=\epsilon_a(j+1/2)$, transport gives
\[
 U_3H_WU_3^*=\epsilon_a(3N_w+B+\tfrac12I),
 \qquad B=\operatorname{diag}(0,1,2).
\]
The diagonal identity preserves the self-adjoint domain defined by
square summability of energy. At $\beta_W\epsilon_a=\log10$, the
quotient has spacing $3\epsilon_a$ and ratio $q$, with the same
thermometer. Its means satisfy $\langle N_w\rangle=1/999$,
$\langle B\rangle=4/37$ and
$3\langle N_w\rangle+\langle B\rangle=1/9$.
Moreover, $s(\rho_r)=s(\rho_q)+s(\sigma_r)$ by exact factorization.

This fixes $g=3\epsilon_a$ when composing \eqref{eq:ii-ter27-lift}
with the Williamson quotient. Preparing $\rho_L\otimes\rho_q$ and
recording projection onto the image of $V$ gives
\[
 V^*(\rho_L\otimes\rho_q)V=\frac{1-q}{Z_L}q^{N+I}.
\]
Its probability is $(1-q)qZ_N/Z_L$, and its conditioned state is
$\rho_N$. The factor $\sigma_r$ is retained through $V\otimes I_3$;
the transported total energy is
$3\epsilon_a(N+1)+\epsilon_a(B+1/2)$.

\subsection{Marked spectral section and complete intensity}
\label{sec:ii-ter27-marcador}

Composition with capacity uses the declared map
$C_{s_n}\mapsto j=C_{s_n}$. By \eqref{eq:ii-ter27-localizacion},
the levels are $23+3n=3(7+n)+2$. This section requires both remainder
$2$ and quotient at least $7$; remainder alone would also admit
$2,5,8,\ldots$. Its projection and probability are
\[
 P_{23}=\sum_{n\geq0}|23+3n\rangle\langle23+3n|,
 \qquad p_{23}=\frac{(1-r)r^{23}}{1-r^3}.
\]
Conditioning and sending $|23+3n\rangle$ to $|n\rangle$ produces
$\rho_q$. The weight $p_{23}$ records the threshold that disappears
from the normalized density.

To incorporate multiplicities $d=(1,380,649)$, define
$\iota|n,a\rangle=|23+3n\rangle\otimes|a\rangle$ in a memory of
dimension $649$, with $1\leq a\leq d_n$. This is an isometry.
Compressing $\rho_r\otimes I_{649}/649$ to its image gives mass
\[
 p_{\mathcal D}=\frac{1-r}{649}r^{23}Z_N(q)
\]
and conditioned state $\rho_N(q)$. If $\lambda_j=(1-r)r^j$, the
intensity relative to the reference population $\lambda_0$ is
\begin{equation}
 B_*:=\frac{\lambda_{23}+380\lambda_{26}+649\lambda_{29}}{\lambda_0}
 =10^{-23}\left(1+\frac{380}{1000}+\frac{649}{1000^2}\right)
 =1.380649\times10^{-23}.
 \label{eq:ii-ter27-intensidad}
\end{equation}
This formula retains spectrum, threshold, multiplicities and reference.
The historical formula was known during construction of this reader;
the operator development specifies its realization and conditions rather
than retrospectively attributing an independent prediction to it.
Interpretation as a map between energy and temperature lines requires
composition of the dimensional readings, treated in the following development.

\subsection{Boundary currents, compression and offset rigidity}
\label{sec:ii-ter27-corrientes}

The area incidence developed in \ref{sec:ii-area-incidencia-explicita}
provides $18+18$ tritic links. Put $u=N_{90}$, $v=N_{120}$, $a=u+v$,
and
\[
 F=3u+4v,\qquad H_R=\epsilon_0F,
 \quad\epsilon_0=30\Delta\tau_R>0,
 \quad\tau_R=\frac{\hbar c}{2\pi R}.
\]
Physical area is $\delta a\,a$, with $\delta a$ fixed by incidence
and Barbero. Provenance remains in $u,v$: equal areas may have different
energies. The currents below use this existing structure while retaining
its Hamiltonian.

On one link, $T^+=|1\rangle\langle0|+|2\rangle\langle1|$ and
$T^-=(T^+)^*$ satisfy $[n,T^\pm]=\pm T^\pm$. Define
\[
 J_{ef}=T^+_{120,f}T^-_{90,e},\qquad
 B_c=\left(\sum_eT^+_{90,e}\right)^3
     \left(\sum_fT^-_{120,f}\right)^2.
\]
The rule $[F,AB]=[F,A]B+A[F,B]$ proves
\[
 [F,J_{ef}]=J_{ef},\quad[a,J_{ef}]=0,
 \qquad[F,B_c]=B_c,\quad[a,B_c]=B_c.
\]
The occupation sequence $(0,2)\to(3,0)\to(2,1)$ has primitive
energies $8,9,10$ and areas $2,3,3$. The same energy increment retains
two distinguishable area actions.

The coefficients $c_j=[z^j](1+z+z^2)^{18}$ count occupations in each
channel: $c_0=1,c_1=18,c_2=171,c_3=1122$. The boundary character
\[
 P(x)=(1+x^3+x^6)^{18}(1+x^4+x^8)^{18}
\]
gives multiplicities $g_7=324,g_8=171,g_9=1122,g_{10}=3078$.
For ranks $d=(1,380,649)$, the first admissible triple of consecutive
energies starts at $8$. Earlier starts are excluded using $g_0=1$,
$g_1=g_2=g_5=0$, $g_3=g_4=18$, $g_6=171$, $g_7=324$ and $g_8=171$:
each earlier triple fails at least one required rank.

A connected compression is constructed by selecting one $(0,2)$ state,
$380$ states $(3,0)$ and $649$ distinct neighbours under the $J_{ef}$.
There are sufficiently many neighbours by a counting proof: every
occupation of sum three has at least two predecessors of sum two,
and each predecessor has at most $18$ successors. The $380$ states
therefore yield at least $\lceil760/18\rceil=43$ distinct predecessors,
and $18\cdot43=774$ neighbours $(2,1)$, from which $649$ may be selected.
The current $B_c$ connects the initial state to all $380$ intermediate
states, and every final neighbour has a selected predecessor. The graph
is connected. Its labels and subspace selection are retained. The based
isometry $\iota_N$ satisfies
\begin{equation}
 H_R\iota_N=\iota_N\epsilon_0(8I+N).
 \label{eq:ii-ter27-compresion}
\end{equation}
Evaluation on each basis vector proves this identity.

Let $P_F=\mathbf1_{N=1}$ and let
$B_N:\mathcal D_0\to\mathcal D_1$, $J_N:\mathcal D_1\to\mathcal D_2$
be the nonzero compressed currents. For
$H_\delta=\epsilon_0(8I+N)+\delta P_F$,
\[
 [H_\delta,B_N]=(\epsilon_0+\delta)B_N,
 \qquad[H_\delta,J_N]=(\epsilon_0-\delta)J_N.
\]
Preserving both increments, or requiring their equality, forces
$\delta=0$. More generally, a diagonal correction preserves each
increment exactly when its values agree at the two ends of every edge.
Connectivity makes it constant. This rigidity classifies diagonal
corrections; its proof does not extend that classification to all
non-diagonal operators.

For $x=e^{-\epsilon_0/\tau}$ and $\tau>0$, Gibbs compression gives
$\iota_N^*e^{-H_R/\tau}\iota_N=x^8x^N$. Its selection probability
from the boundary is $x^8Z_N(x)/P(x)$. Recording selection and complement
through $P_N\otimes|s\rangle+(I-P_N)\otimes|f\rangle$ is isometric
and preserves energy for degenerate flags. On each upward edge, rates
$xr_e$ and reverse rate $r_e>0$ satisfy detailed balance for weights
$x^N$. Connectivity proves uniqueness of the stationary state; grade
masses are $d_nx^n/Z_N(x)$. This specifies a thermal preparation with
a bath rather than an isolated unitary evolution.

\Needspace{29\baselineskip}
\subsection{Thermal equivalence with transported reference and scale}
\label{sec:ii-ter27-equivalencia}

\begin{theorem}[Transport of the centred thermal exponent]
\label{thm:ii-ter27-transporte}
Let $H_1,H_2$ be finite self-adjoint operators, $\tau_i>0$, and $T$
an isometry whose image reduces $H_2$. Traces in the second space are
taken on that image. Let $\sigma_{\mathrm{ref}}$ be a normalized state
on the first space, and define the transported reference energies by
\[
 e_1=\operatorname{Tr}(\sigma_{\mathrm{ref}}H_1),\qquad
 e_2=\operatorname{Tr}(T\sigma_{\mathrm{ref}}T^*H_2).
\]
If
\[
 \frac{H_2-e_2I}{\tau_2}T
 =T\frac{H_1-e_1I}{\tau_1},
\]
then the centred partition function, Gibbs state, entropy and probabilities
of all transported observables, including the marker, are preserved.
Conversely, for faithful states, equality of the transported states
implies this identity with the references defined above.
\end{theorem}
\begin{proof}
Functional calculus intertwines exponentials of both exponents. Trace
invariance under an isometry onto its image gives equality of centred
partitions; normalization yields $\rho_2=T\rho_1T^*$. Trace cyclicity
transports observables, and spectral calculus transports entropy.
Conversely, logarithms of faithful states give equality of $H_i/\tau_i$
up to a scalar, the difference of their log partition functions.
Taking its expectation in $\sigma_{\mathrm{ref}}$ identifies that scalar
with $e_2/\tau_2-e_1/\tau_1$; subtracting it gives the stated centred identity.
\end{proof}

For the capacity reader, $C=23I+3N$ and
\begin{equation}
 H_{\rm cap}=\frac{\epsilon_a\log10}{\log3}C,
 \quad\tau_{\rm clk}=\frac{\epsilon_a}{\log3},\qquad
 H_{W,\rm rel}=\epsilon_a C,
 \quad\tau_W=\frac{\epsilon_a}{\log10}.
 \label{eq:ii-ter27-hcap}
\end{equation}
Both exponents equal $C\log10$ although energies and temperatures have
different scales. On the boundary compression, preparation at
$\tau_*=\epsilon_0/\log1000$ satisfies
\[
 \frac{H_R-8\epsilon_0I}{\tau_*}\iota_N
 =\iota_N\frac{H_{\rm cap}-23\tau_{\rm clk}\log10\,I}{\tau_{\rm clk}}
 =\iota_N N\log1000.
\]
The native horizon activity remains $e^{-30\Delta}$; this preparation
has ratio $\tau_*/\tau_R=30\Delta/\log1000$. Varying $R$ changes
$\epsilon_0$ and $\tau_R$ together, preserving their quotient.
Before centring, boundary and capacity exponents differ by
$\log10\,I$, because $8\log1000-23\log10=\log10$. Their traces are
$10^{-24}Z_N$ and $10^{-23}Z_N$. The factor $1/10$ also occurs in
their references and must be retained when comparing absolute intensities.

If $\iota_1,\iota_2$ realize the same ranks, grading and marker,
$U=\iota_2\iota_1^*$ is unitary between their images and intertwines
Hamiltonian and marker. The preceding theorem implies the same reader,
all its spectral moments and its thermal derivatives for any admissible
basis choice. A unique microscopic representative is unnecessary for
that conclusion. Comparing area, currents or relaxation times also
requires transport of those additional observables.

Finally, for
$P_0=|\Omega_{\rm ext}\rangle\langle\Omega_{\rm ext}|$ and
$\rho_N=q^N/Z_N$, one has
\[
 p_0=\operatorname{Tr}(P_0\rho_N)=Z_N^{-1},\qquad
 Z_N=1+\frac{\operatorname{Var}_{\rho_N}(P_0)}{p_0^2}.
\]
The second identity uses $P_0^2=P_0$, hence variance $p_0(1-p_0)$.
Normalized state and marker recover the relative trace; the recorded
threshold recovers the factor in \eqref{eq:ii-ter27-intensidad}.
The result thereby preserves the structure, preparation, reference and
memory required for composition with the action--temperature reading.
